\documentclass[12pt,a4paper]{article} %12 кегль, формат бумаги А4, тип статья

\usepackage[T2A]{fontenc} %поддержка кириллицы
\usepackage[utf8]{inputenc} %кодировка символов юникод
\usepackage[english,russian]{babel} %список рабочих языков
\usepackage{cmap} %для поиска в pdf
\usepackage[left=2cm,right=2cm,top=2cm,bottom=2cm]{geometry} %размеры отступов от краев страницы
\usepackage{graphicx} %возможность работать с графикой
\usepackage{enumitem} %для работы со списками
\usepackage{float} %для точного размещения таблиц и графики
\usepackage{caption} %для создания подписей таблиц и графики
\usepackage{amsmath,amssymb,amsthm,amsfonts} %штуки для математики
\usepackage{extarrows} %чтобы подписи над и под нестандартными стрелками писать короче в коде
\usepackage{makeidx} %собирает ключевые слова для предметного указателя
\usepackage[hidelinks]{hyperref} %убирает выделение ссылок; добавляет гиперссылки в оглавление, перекрестные ссылки и URL
\usepackage{setspace} %устанавливает отступы
\linespread{1.5} %длина межстрочного интервала
\vspace{125mm} %длина абзацного отступа
\usepackage{titlesec}
\titleformat{\section}{\normalfont\Large\bfseries\centering}{}{0pt}{} % эти две строчки убирают циферки перед названиями разделов в основном тексте и выравнивают названия секций по центру
\usepackage{paratype} %чтобы точно был поиск и шрифт хороший

\author{Автор конспекта: Краснов Борис} %автор
\title{Белов Ю.С.: Целые функции} %название
\date{Весна 2026} %дата создания документа

\newtheorem{theorem}{Теорема}
\newtheorem{lemma}{Лемма}
\newtheorem{definition}{Определение}
\newtheorem{proposition}{Предложение}
\newtheorem{consequence}{Следствие}
\newcommand{\R}{\mathbb{R}}
\newcommand{\Co}{\mathbb{C}}
\newcommand{\E}{\mathbb{E}}
\newcommand{\D}{\mathbb{D}}
\newcommand{\PR}{\mathbb{P}}
\newcommand{\charf}[2][t]{\varphi_{#2}(#1)}
\newcommand{\F}{\mathcal{F}}
\newcommand{\M}{\mathfrak{M}}